\documentclass[10pt,a4paper]{amsart}
\usepackage[margin=19mm]{geometry}
\usepackage{amsmath,amssymb,mathtools}
\usepackage[scr=rsfs,cal=euler]{mathalfa}
\usepackage{xfrac}
\usepackage[unicode,hidelinks,bookmarks=false]{hyperref}
\newcommand{\ie}{i.e.\ }
\newcommand{\cf}{cf.\ }
\newcommand{\resp}{respectively\ }
\newcommand{\Subsection}{\subsection}
\providecommand{\nobreakdash}{\nobreak\hbox{-}}
\setlength{\emergencystretch}{2em}
\multlinegap0pt
\usepackage{amssymb}
\usepackage{xcolor}
\usepackage{algorithm}\usepackage{algpseudocode}
\newtheorem{lemma}{Lemma}
\newtheorem{definition}{Definition}
\newtheorem{theorem}{Theorem}
\title{Article Title}
\author{Soumic Sarkar}
\date{Source: arXiv:2608.07552v1 (CC BY 4.0)}
\begin{document}
\maketitle

\noindent\emph{Source and licence.} Article Title. Version arXiv:2608.07552v1 (CC BY 4.0), distributed by arXiv under the Creative Commons Attribution 4.0 International licence, http://creativecommons.org/licenses/by/4.0/, as recorded in the arXiv licence metadata for this exact version and shown on the abstract page https://arxiv.org/abs/2608.07552v1. Full licence notice: \texttt{licenses/arxiv-2608.07552-CC-BY-4.0.txt}.\par\smallskip

\noindent\textit{Editorial scope.} Counted unit: the article's theorem thm:global (source0.tex lines 267--270). The article's complete original proof of it is delivered with it (source0.tex lines 272--274, one \texttt{proof} environment of 3 lines). No mathematics is written, repaired or re-derived: the statement and the proof are the article's own lines, in the article's own notation, and no retained line is substituted or re-flowed.  Retained uncounted as the source-local context the proof needs: lines 255--258; lines 262--265.  Nothing was omitted from any retained span, so the package carries no omission record.  No retained line contains a citation, so the wrapper carries no bibliography and no bibliography entry.  No retained line contains a figure environment, an image-inclusion command or drawing code, so no image is delivered and the package declares no asset dependency.  Editorial formatting only, in addition: an AMS \texttt{amsart} preamble that restates the article's own declarations and macros used by the retained lines, namely the environments \texttt{lemma}, \texttt{definition}, \texttt{theorem} on separate counters, together with the standard packages those lines need.  The retained environments print in this document as Lemma 1, Definition 1, Theorem 1 in order of appearance, whereas the article prints the same items with the numbering of its own full document.  The complete native source of the article as distributed in the arXiv e-print for this exact version is bundled unchanged as \texttt{sources/207092/source0.tex}.
\begin{lemma}[Dini comparison]
\label{lem:dini}
Let $\psi:[0,\infty)\to\mathbb{R}_{\ge0}$ be continuous and suppose $D^+\psi(t)\le-\lambda\psi(t)$ for all $t\ge0$ and some $\lambda>0$, where $D^+$ denotes the upper Dini derivative $D^+\psi(t)=\limsup_{\epsilon\to0^+}(\psi(t+\epsilon)-\psi(t))/\epsilon$. Then $\psi(t)\le e^{-\lambda t}\psi(0)$ for all $t\ge0$.
\end{lemma}

\begin{definition}
\label{def:intrinsiccontraction}
System $\dot z=f(z)$ is intrinsically contracting with rate $\lambda>0$ if $D^+F_K(z(t),\delta z(t))\le-\lambda F_K(z(t),\delta z(t))$ holds along every trajectory and every admissible infinitesimal variation.
\end{definition}

\begin{theorem}
\label{thm:global}
If $\dot z=f(z)$ is intrinsically contracting with rate $\lambda>0$, then $d_K(\phi_t(x),\phi_t(y))\le e^{-\lambda t}d_K(x,y)$ for all $x,y\in D$, $t\ge0$.
\end{theorem}

\begin{proof}
Fix $x,y\in D$ and a piecewise smooth curve $\gamma$ joining them, and write $\gamma_t=\phi_t\circ\gamma$. For each fixed $s$, $t\mapsto F_K(\gamma_t(s),\dot\gamma_t(s))$ is continuous, and Definition~\ref{def:intrinsiccontraction} gives the hypothesis of Lemma~\ref{lem:dini} for $\psi(t)=F_K(\gamma_t(s),\dot\gamma_t(s))$, yielding $F_K(\gamma_t(s),\dot\gamma_t(s))\le e^{-\lambda t}F_K(\gamma(s),\dot\gamma(s))$ for every $s$. Integrating over $s\in[0,1]$ and taking the infimum over $\gamma$ on both sides gives the stated bound.
\end{proof}

\end{document}
